\documentclass[aspectratio=169, xcolor=table]{beamer}

% Theme (header progress can be disabled: [headprogress=false])
\usepackage[headprogress=true]{beamerthemeTUDneo}

% Math, tables, figures, code, algorithms
\usepackage{amsmath, amssymb, mathtools, bm}
\usepackage{booktabs, array}
\usepackage{siunitx}
\usepackage{graphicx}
\usepackage{caption}
\usepackage{subcaption}
\usepackage{listings}
\usepackage{algorithm}
\usepackage{algpseudocode}

% Bibliography (biber)
\usepackage[backend=biber,style=numeric,sorting=none]{biblatex}
\addbibresource{../ref/refs.bib}

% Listings style (simple, neutral)
\lstset{
  basicstyle=\ttfamily\small,
  frame=single,
  breaklines=true,
  numbers=none,
  keywordstyle=\color{TUDgray}\bfseries,
  commentstyle=\color{TUDgray},
  rulecolor=\color{black}
}

% Show a compact Outline at each (sub)section start
\AtBeginSection{
  \begin{frame}
    \frametitle{Outline}
    \tableofcontents[currentsection,hideallsubsections]
  \end{frame}
}
\AtBeginSubsection{
  \begin{frame}
    \frametitle{Outline}
    \tableofcontents[currentsubsection]
  \end{frame}
}

% Title info
\title[TUD Neo Demo]{TUD Neo — Modern White Theme}
\subtitle{Clean footline, geometric accent}
\author[Sihao Wei]{Sihao Wei}
\institute[Faculty of Transportation and Traffic Sciences]{Faculty of Transportation and Traffic Sciences}
\date{\today}

\begin{document}

% Title frame (white center, centered logo)
\begin{frame}[plain]
  \vspace*{1.2cm}
  \centering
  {\usebeamerfont{title}\usebeamercolor[fg]{structure}\LARGE \inserttitle\par}
  \vspace{0.4cm}
  {\normalsize \insertsubtitle\par}
  \vspace{0.8cm}
  {\large \insertauthor\par}
  {\small \insertinstitute\par}
  \vspace{0.6cm}
  {\small \insertdate\par}
  \vspace{0.8cm}
  % Place your SVG at assets/Logo_TU_Dresden_2025.svg (or PDF/PNG/JPG)
  \TUDCenteredLogo[0.16\paperwidth]{../assets/Logo_TU_Dresden_2025}
\end{frame}

% ---------------- Outline ----------------
\begin{frame}{Outline}
  \tableofcontents
\end{frame}

\section{Bullets \& Math}
\subsection{Inline and Display}

\begin{frame}{Bullets \& Math}
  \textbf{Item:} Clean white background with TU blue accents inspired by TU Dresden's visual identity~\cite{dantzig1951}.\\
  \textbf{Note:} Readable footline; total pages + pages left in the header~\cite{berge1963}.

  \medskip
  Inline math: \( f(x)=x^2 \), bold vector \( \bm{v}\in\mathbb{R}^3 \).

  \medskip
  Displayed integral:
  \[
    \int_0^1 x^m\,dx = \frac{1}{m+1},\quad m> -1.
  \]
\end{frame}

\subsection{Blocks}

\begin{frame}{Blocks}
  \begin{block}{Strong statement}
    Clean block environments; titles use TU blue with white text.
  \end{block}
  \begin{exampleblock}{Example}
    Use white titles with blue text for helpful tips.
  \end{exampleblock}
\end{frame}

\section{Algorithms}
\subsection{Pseudocode}
\begin{frame}{Pseudocode}
  \begin{algorithm}[H]
    \caption{Greedy Selection (toy)}
    \begin{algorithmic}[1]
      \Require set $S$, score $w(\cdot)$, budget $k$
      \Ensure  subset $A$
      \State $A \gets \emptyset$
      \While{$|A| < k$}
        \State $x \gets \arg\max_{y\in S\setminus A} w(y)$
        \State $A \gets A \cup \{x\}$
      \EndWhile
      \State \Return $A$
    \end{algorithmic}
  \end{algorithm}
\end{frame}

\section{Tables \& Figures}
\subsection{Table}
\begin{frame}{Table}
  \begin{table}
    \centering
    \caption{Tiny example with \texttt{booktabs}.}
    \begin{tabular}{@{}lS[table-format=1.2]S[table-format=2.1]@{}}
      \toprule
      Method & {Accuracy} & {Time (s)} \\
      \midrule
      A & 0.92 & 12.3 \\
      B & 0.89 & 9.8 \\
      \bottomrule
    \end{tabular}
  \end{table}
\end{frame}

\subsection{Figure}
\begin{frame}{Figure}
  \begin{figure}
    \centering
    % Put a sample figure at assets/example-graph.(pdf|png|jpg)
    \includegraphics[width=.62\linewidth]{../assets/example-graph}
    \caption{Example figure (use \texttt{example-graph})~\cite{clarke1964}.}
  \end{figure}
\end{frame}

\section{Code}
\begin{frame}[fragile]{Listings}
\begin{lstlisting}[language=Python,caption={Simple Python snippet}]
def f(x: float) -> float:
    return x*x + 1.0

print(f(2.0))
\end{lstlisting}
\end{frame}

\section{References}
\begin{frame}[allowframebreaks]{Bibliography}
  \footnotesize
  \begin{enumerate}
    \item \fullcite{dantzig1951}
    \item \fullcite{berge1963}
    \item \fullcite{clarke1964}
  \end{enumerate}
\end{frame}

\end{document}
